\documentclass[11pt]{article}
\usepackage[margin=1in]{geometry}
\usepackage{amsmath,amssymb,amsthm}
\usepackage[T1]{fontenc}
\usepackage{lmodern}
\usepackage{graphicx}
\usepackage{hyperref}

\newtheorem{theorem}{Theorem}
\newtheorem{lemma}{Lemma}
\newtheorem{corollary}{Corollary}

\newcommand{\eps}{\varepsilon}
\newcommand{\R}{\mathbb R}
\newcommand{\cL}{\mathcal L}
\newcommand{\Ext}{\operatorname{Ext}}
\newcommand{\perpB}{\perp_B}
\newcommand{\perpBe}{\perp_B^\varepsilon}

\title{Property P for two-dimensional polyhedral Banach spaces}
\author{Codex agent lane 16}
\date{July 19, 2026}

\begin{document}
\maketitle

\section*{Source Question}

Mandal--Manna--Paul--Sain define Property P for a Banach space \(X\):
there should exist \(\eps\in(0,1)\) such that every bounded linear operator
\(T\in\cL(X)\) preserving \(\eps\)-approximate Birkhoff--James orthogonality at
every point is a scalar multiple of an isometry \cite{MMPS}.  Their source page
contains the question and definition:

\begin{center}
\includegraphics[width=0.88\linewidth]{figures/open_problem_crop.png}
\end{center}

The source proves that regular \(2n\)-gon planes have Property P for \(n>2\),
and observes that the regular \(4\)-gon case fails because it is isometric to
\(\ell_1^2\).  This packet removes the regularity assumption and gives the full
two-dimensional polyhedral classification.

\section*{Source Tools Used}

We use three facts from \cite{MMPS}.  First, if a nonzero operator preserves
\(\eps\)-orthogonality for some \(0\leq\eps<1\), then it is injective
\cite[Lemma 2.7]{MMPS}.  Second, for each finite-dimensional polyhedral space
\(X\) the source defines a positive constant \(\eps_X\) such that, for
\(0\leq\eps<\eps_X\), an \(\eps\)-orthogonality preserver \(T\in\cL(X)\) maps
nonsmooth points of \(X\) to nonsmooth points of \(X\)
\cite[Corollary 2.16]{MMPS}.  In a polyhedral plane, the nonsmooth boundary
points are exactly the vertex rays.  Third, the exact
Blanco--Koldobsky--Turnsek theorem says that a linear operator preserving exact
Birkhoff--James orthogonality at every point is a scalar multiple of an isometry
\cite{Koldobsky,BlancoTurnsek}.

\section*{Proof Intuition}

For a polygonal unit ball, sufficiently small approximate preservers are already
forced by the source's lemmas to respect the finite set of vertex rays.  If the
polygon has at least six vertices, then it has at least three projective vertex
lines.  A linear map that sends these finitely many lines to themselves, in some
order, is determined up to scalar once the induced line map is fixed; there are
only finitely many such possibilities after normalization.

If Property P failed, we could find normalized nonsimilarity preservers with
\(\eps_k\downarrow0\).  Passing to a subsequence, the finite vertex-ray pattern
would stabilize and the operators would converge to an exact orthogonality
preserver.  The exact Blanco--Koldobsky--Turnsek theorem makes the limit a
similarity.  But inside a fixed vertex-ray pattern there is only one normalized
operator up to sign/scalar, so the whole stabilized sequence was already made of
similarities.  This contradiction proves Property P.

\begin{lemma}[Finite ray rigidity]
Let \(X\) be two-dimensional and let \(\mathcal R\) be a finite set of at least
three one-dimensional subspaces of \(X\).  Then the set of invertible operators
which send every line in \(\mathcal R\) into a line in \(\mathcal R\) is a finite
union of one-dimensional subspaces of \(\cL(X)\).
\end{lemma}

\begin{proof}
For each map \(\sigma:\mathcal R\to\mathcal R\), let
\[
  M_\sigma=\{A\in\cL(X): A(L)\subseteq \sigma(L)\text{ for every }L\in\mathcal R\}.
\]
There are only finitely many \(\sigma\)'s, so it is enough to show that whenever
\(M_\sigma\) contains an invertible operator, \(M_\sigma\) is at most
one-dimensional.

Choose three distinct lines \(L_1,L_2,L_3\in\mathcal R\), with \(L_1,L_2\)
independent.  Since \(M_\sigma\) contains an invertible operator,
\(\sigma(L_1)\) and \(\sigma(L_2)\) are independent.  Pick nonzero
\(e_i\in L_i\) and nonzero \(u_i\in\sigma(L_i)\) for \(i=1,2\).  After scaling
\(e_1,e_2\), write \(L_3=\R(e_1+e_2)\).  For \(A\in M_\sigma\), there are
scalars \(a,b\) such that \(Ae_1=a u_1\) and \(Ae_2=b u_2\).  The extra
condition \(A(L_3)\subseteq\sigma(L_3)\) says that \(a u_1+b u_2\) lies in a
fixed line.  Because \(u_1,u_2\) are independent, this imposes a fixed linear
relation between \(a\) and \(b\).  Thus the possible pairs \((a,b)\), and hence
the possible operators \(A\), form a space of dimension at most one.
\end{proof}

\begin{lemma}[Positive gap for nonsimilar ray maps]
Let \(X\) be a two-dimensional polyhedral Banach space whose unit sphere has at
least six vertices.  Among norm-one operators that map every vertex line of
\(B_X\) into a vertex line, the nonsimilarities have a positive lower bound on
the \(\eps\) needed to preserve \(\eps\)-orthogonality at every point.
\end{lemma}

\begin{proof}
Let \(\mathcal R\) be the finite set of vertex lines.  By the previous lemma,
the norm-one ray-preserving operators lie in a finite set, up to sign.  Let
\(\mathcal F\) be representatives of those which are not scalar multiples of
isometries.

For \(A\in\mathcal F\), exact Blanco--Koldobsky--Turnsek implies that \(A\) does
not preserve exact Birkhoff--James orthogonality.  Hence there are unit vectors
\(x,y\in X\) with \(x\perp_B y\) but \(Ax\not\perp_B Ay\).  The
Chmielinski--Stypula--Wojcik characterization used in the source gives
\[
  \min_{f\in J(Ax)}\frac{|f(Ay)|}{\|Ay\|}>0,
\]
because \(J(Ax)\) is compact and no member of \(J(Ax)\) annihilates \(Ay\).
Therefore \(A\) cannot preserve \(\eta_A\)-orthogonality for all
\(\eta_A\) below this positive number.  Since \(\mathcal F\) is finite, the
minimum of these positive lower bounds is again positive.
\end{proof}

\begin{theorem}[Classification in polyhedral planes]
Let \(X\) be a two-dimensional real polyhedral Banach space.  Then \(X\) has
Property P if and only if the unit sphere of \(X\) has at least six extreme
points.
\end{theorem}

\begin{proof}
First suppose that \(S_X\) has four extreme points.  Then \(B_X\) is a
centrally symmetric quadrilateral, hence \(X\) is linearly isometric to
\(\ell_\infty^2\), equivalently to \(\ell_1^2\).  The source proves that
\(\ell_p^n\) fails Property P for \(1\leq p<\infty\); in particular
\(\ell_1^2\) fails Property P \cite[Corollary 2.27]{MMPS}.  Thus \(X\) fails
Property P.

Conversely suppose that \(S_X\) has \(2m\) extreme points with \(m\geq3\).
Let \(\mathcal R\) be the set of vertex lines.  Let \(\delta>0\) be the positive
gap from the previous lemma, and choose
\[
  0<\eps_0<\min\{\eps_X,\delta\}.
\]
We claim that \(\eps_0\) witnesses Property P.

Let \(T\in\cL(X)\) preserve \(\eps_0\)-orthogonality at every point.  If
\(T=0\), then \(T\) is a scalar multiple of an isometry.  If \(T\neq0\), then
the source injectivity lemma makes \(T\) invertible.  Since
\(\eps_0<\eps_X\), the source nonsmooth-point preservation corollary sends each
vertex ray into a vertex ray.  Thus \(\|T\|^{-1}T\) is a norm-one
ray-preserving operator.  If it were not a scalar multiple of an isometry, the
positive gap lemma would force the required preservation constant to be at least
\(\delta>\eps_0\), contradiction.  Hence \(\|T\|^{-1}T\), and therefore \(T\),
is a scalar multiple of an isometry.
\end{proof}

\section*{Verification and Novelty Check}

This result is a partial result for the global Property P question, but it is a
full classification inside two-dimensional polyhedral spaces.  It strictly
extends the source's regular polygon theorem: regularity is not used.

Bounded duplicate/literature check performed before promotion:
\begin{itemize}
  \item local run indexes were searched for arXiv:2501.10719, Property P,
  approximate preservation of Birkhoff--James orthogonality, and two-dimensional
  polyhedral variants;
  \item web search located the source's published Colloquium Mathematicum entry
  and the later Lebesgue--Bochner continuation arXiv:2601.02786;
  \item targeted web searches for arbitrary polygon or two-dimensional
  polyhedral Property P classifications found no exact statement beyond the
  source's regular polygon theorem.
\end{itemize}

Human-review recommendation: verify that the source constant \(\eps_X\) is
strictly positive for every two-dimensional polyhedral space and that
\cite[Corollary 2.16]{MMPS} is applied exactly as stated.  The finite-ray
compactness argument is elementary once those source lemmas are accepted.

\begin{thebibliography}{9}

\bibitem{MMPS}
K. Mandal, J. Manna, K. Paul, and D. Sain,
\emph{On approximate preservation of orthogonality and its application to
isometries}, arXiv:2501.10719; Colloquium Mathematicum 179 (2025), 189--211.

\bibitem{Koldobsky}
A. Koldobsky,
\emph{Operators preserving orthogonality are isometries}, Proc. Roy. Soc.
Edinburgh Sect. A 123 (1993), 835--837.

\bibitem{BlancoTurnsek}
A. Blanco and A. Turnsek,
\emph{On maps that preserve orthogonality in normed spaces}, Proc. Roy. Soc.
Edinburgh Sect. A 136 (2006), 709--716.

\bibitem{CSW}
J. Chmielinski, T. Stypula, and P. Wojcik,
\emph{Approximate orthogonality in normed linear spaces and its applications},
Linear Algebra Appl. 531 (2017), 305--317.

\bibitem{MJ}
Mohit and R. Jain,
\emph{Approximate Birkhoff-James orthogonality preserver on
Lebesgue-Bochner spaces}, arXiv:2601.02786, 2026.

\end{thebibliography}

\end{document}
